\section{Agent Swarm y patrones de colaboración}

\subsection{Niveles de la estructura de Swarm}

El Swarm que menciona Omar generalmente incluye three tiers: 1) CLI/Coordinator, 2) execution agents, 3) verification agents. El Coordinator se encarga de la orchestration, los execution agents procesan las prompt/LLM calls, los verification agents revisan los outputs y asignan un `evidence rating`.

\begin{importantbox}{El valor de la separación de niveles}
La separación de niveles reduce la complejidad. Incluso si un execution agent comete un error en su output, el verification agent puede bloquear su paso por el pipeline en el nivel de la `Evidence Matrix`.
\end{importantbox}

\subsection{Colaboración de la cadena de herramientas}

Los tools del Swarm generalmente se exponen en forma de `MCP`; DSPy registra un `ToolDescriptor` al definir un Module. En runtime, el Coordinator invoca los tools en el orden que indica la `signature`, y registra los parámetros de la llamada y el valor de retorno en el evidence log.

\begin{itemize}
    \item `search\_tool`: invoca rápidamente el retrieval index
    \item `cohere\_tool`: invoca el LLM for consistent reasoning
    \item `validate\_tool`: run schema + rule-based checks sobre el tool output
\end{itemize}

\subsection{Resumen de la sección}

El éxito del Agent Swarm depende de la estructura jerárquica y la coordinación de herramientas: solo cuando el Coordinator, el Validator y el Executor cumplen cada uno su función es posible mantener bajo control el compound system.

